% Frozen main-section module, 29 October 2026.
% Requires amsmath, amssymb, amsthm; theorem, lemma, proposition environments.
% General derivative appendix supplies Theorem thm:mc-derivative.
\section{A simultaneous-cut improvement of the pentagon product}
\label{sec:hko-ascent}

The pentagon product is not a local maximum when the number of facets is
unrestricted. We prove this by a fixed family of twenty supporting
inequalities. The calculation combines the simultaneous-splitting derivative
with a reduction tailored to its cut profiles and an exhaustive exact
calculation of the remaining shortest-word responses. No particular nonzero
improving parameter is certified here.

Let $P\subset\mathbb R^2$ be the regular pentagon with vertices
$(\cos(2\pi k/5),\sin(2\pi k/5))$, $0\leq k<5$, and let
$J_2(x,y)=(-y,x)$. In coordinates $(q,p)$ put
\[
 H=P\times_L(-J_2P),\qquad
 J(q,p)=(-p,q),\qquad \tau=\tan(\pi/5).
\]
Let $n_i$ be the support-normalized outward normal of the edge from vertex
$i$ to vertex $i+1$ of $P$, with indices taken modulo five, and set
$m_i=-J_2n_i$. Thus $n_0=(1,\tau)$ and
$|n_i|=|m_i|=\sqrt5-1$. The ten normals of $H$ are
$a_i=(n_i,0)$ and $a_{5+i}=(0,m_i)$. Define
\begin{align}
 D_i&=\left(-\frac35n_i,\frac{J_2n_i}{|n_i|}\right),&
 D_{5+i}&=\left(-\frac{J_2m_i}{|m_i|},-\frac35m_i\right),
 \label{eq:ha-directions}\\
 H_t&=\{x:a_j\cdot x\leq1,\ (a_j+tD_j)\cdot x\leq1,
                         \ 0\leq j<10\}.\label{eq:ha-family}
\end{align}
Every old inequality is retained, so $H_t\subset H$ for $t\geq0$.
For sufficiently small $t$ these are full-dimensional bounded polytopes and
$H_t\to H$ in Hausdorff distance. Equivalently,
\[
 H_t=H\cap L_t^{-1}H,\qquad
 L_t=\begin{pmatrix}aI_2&-bJ_2\\ bJ_2&aI_2\end{pmatrix},\quad
 a=1-\frac35t,\quad b=\cos(\pi/5)t.
\]
There are exactly twenty genuine facets for all sufficiently small $t>0$.
Indeed, on the relative interior of each old facet, $D_i\cdot x$ takes
both positive and negative values on open patches. A negative patch retains
the old facet with every competing inequality strict. At a positive patch,
the point $x/(1+tD_i\cdot x)$ lies on the new hyperplane, strictly satisfies
the old inequality, and keeps all other inequalities strict for small $t$.
Thus each old and each new row supports a facet.

\begin{theorem}[Strict inward ascent]\label{thm:ha-ascent}
For the family~\eqref{eq:ha-family}, with
$\operatorname{sys}(K)=c_{\mathrm{EHZ}}(K)^2/(2\operatorname{Vol}_4(K))$,
\begin{equation}\label{eq:ha-slope}
 \left.\frac{d}{dt}\right|_{0+}\log\operatorname{sys}(H_t)
       =\frac{20+64\sqrt5}{9375}>0.
\end{equation}
Consequently $\operatorname{sys}(H_t)>\operatorname{sys}(H)$ for every
sufficiently small positive $t$. In particular, $H$ is not a local maximum
among convex bodies in the Hausdorff topology.
\end{theorem}

The statement concerns at most twenty facets and does not contradict local
maximality with an upper bound of eleven facets. The existence of a small
improving interval follows from the positive derivative; the proof gives no
explicit endpoint for that interval.

\subsection{The cut profile and the repeated-word issue}

On each original facet the inward profile is
\begin{equation}\label{eq:ha-profile}
 f_i=(\ell_i-3/5)_+,\qquad
 \ell_i(q,p)=\frac{J_2n_i}{|n_i|}\cdot p,
 \qquad
 f_{5+i}=(\ell_{5+i}-3/5)_+,\quad
 \ell_{5+i}(q,p)=-\frac{J_2m_i}{|m_i|}\cdot q.
\end{equation}
These profiles are convex on their respective facets and constant along
their own characteristic directions $Ja_i$. The constancy is essential:
general convex profiles need not admit the reduction below.

Use the active-word action
$Q(\sigma,\beta)=\sum_{r<s}\beta_r\beta_s
\langle Ja_{\sigma_r},a_{\sigma_s}\rangle$, with
$\sum\beta_ra_{\sigma_r}=0$, $\sum\beta_r=1$, and $\beta_r\geq0$.
Write $Q_*=1/(2c_{\mathrm{EHZ}}(H))=\tau/5$.
For a maximizing word, $\mu_\gamma$ denotes its normalized-time occupation
measure on the disjoint union of labeled facets. The general derivative
theorem uses all finite maximizing words, including repeated labels;
an existence theorem for a simple shortest word alone would not suffice.

\begin{lemma}[Removal of repeated visits]\label{lem:ha-repeat}
Let $K$ be a Lagrangian product of planar polygons with no antipodal facet
normals. Suppose each continuous convex facet profile is constant on lines
parallel to its own characteristic direction. Every finite maximizing word
has a simple maximizing word with at least as large an occupation response.
The same supremum therefore bounds the weak closure of the finite-word
occupation measures.
\end{lemma}

\begin{proof}
Delete zero weights. Suppose a full normal $a=(n,0)$ occurs twice, at
positions $r<s$. Transfer weight by
$\beta_r(z)=\beta_r+z$, $\beta_s(z)=\beta_s-z$ on
$[-\beta_r,\beta_s]$, fixing all other weights. Closure and normalization
are preserved. The action is affine in $z$, since its quadratic coefficient
pairs $a$ with itself. The repeated-word polar-loop bound $Q\leq Q_*$,
together with the interior maximum at $z=0$, makes this affine function
constant. Equality of the normals alone would not imply constancy.

At every interior parameter, the active normals span both factor planes:
a positive closed combination without antipodal directions contains at least
three distinct directions. The augmented normal matrix has rank five,
because averaging a relation $u\cdot a_j+c=0$ against the positive weights
first gives $c=0$, then $u=0$. Thus the KKT multiplier $\lambda(z)$ is unique
and affine in $z$. With $T=(2Q_*)^{-1}$, the reconstructed vertices are
\[
 x_k(z)=-T\lambda(z)+2TJ\sum_{j<k}\beta_j(z)a_{\sigma_j}.
\]
They lie in $K$ on their labeled facets by the reconstruction argument in
Theorem~\ref{thm:mc-derivative}.

Choose the cyclic start just before the first duplicate and put $v=2TJa$.
Every fixed-weight segment translates with derivative $u$ or $u+v$,
according to which interval between the duplicates contains it. All
q-facet equations give $n_j\cdot u_q=0$, so $u_q=0$. The p-facet equations
give $m_j\cdot u_p=0$ outside the interval and
$m_j\cdot(u_p+v_p)=0$ inside. There are at least three distinct p-normal
directions; one interval contains two nonparallel ones. Hence $u_p=0$ or
$u_p=-v_p$, and therefore $u=0$ or $u=-v$.

Both duplicate segments consequently move and change length only along
their own characteristic direction. Their profile values $c_r,c_s$ stay
constant, and their combined response
$(\beta_r+z)c_r+(\beta_s-z)c_s$ is affine. Every other weight is fixed and
its segment translates affinely; integrating its convex profile gives a
convex function of $z$. The total response is convex. By continuity it
extends to the endpoints, where the zero-weight occurrence is deleted.
One endpoint has at least the original response. The p-normal case is
identical with the factors exchanged. Iteration removes all repetitions.
Continuity on the compact disjoint facet union passes the bound to weak
limits of the occupation measures.
\end{proof}

\subsection{Complete calculation of the simple-word response}

\begin{proposition}[Exact finite response certificate]\label{prop:ha-certificate}
For the profiles~\eqref{eq:ha-profile},
\begin{equation}\label{eq:ha-response}
 \max_{\text{simple maximizing }\gamma}\int f\,d\mu_\gamma=\frac1{125}.
\end{equation}
The complete maximizing weight family is covered by 190 charts: 50 points
on six-label words and 140 closed intervals on seven-label words.
\end{proposition}

\begin{proof}
We describe both the exhaustive reduction and the exact predicates checked
by the accompanying portable certificate.
The normalized closure polytope in either factor is
\[
 C_q=\{u\geq0:\sum u_i=1,\ \sum u_in_i=0\},\qquad
 C_p=\{v\geq0:\sum v_i=1,\ \sum v_im_i=0\}.
\]
Its vertices have support at most three by rank of the three affine
constraints. No support of size one or two closes. Solving all ten
three-element supports gives precisely five positive vertices in each
factor, with exact weights in $\mathbb Q(\tau)$.

For a fixed word, split its weights as $(tu,(1-t)v)$. Same-factor pairings
vanish, so $Q=t(1-t)B_\sigma(u,v)$ with $B_\sigma$ bilinear.
Every positive global maximizer has $t=1/2$. Decomposing $u,v$ into closure
vertices expresses its value as a convex combination of triangle-pair
actions. Each pair action is at most $Q_*$, so every pair with positive
product coefficient must attain $Q_*$. Conversely, mixtures supported on
a complete bipartite set of winning pairs attain $Q_*$.

There are $25\cdot5!=3000$ triangle-pair cyclic six-orders. Exact evaluation
checks that all have action at most $\tau/5$, with precisely 50 equality
words. Every simple support word extends by zero weights to a cyclic
ten-order. Fixing label zero first gives all $9!=362880$ such orders.
For each, form the $5\times5$ graph of its winning induced triangle pairs.
Every maximizing weight pair is covered by a maximal complete bipartite
subgraph. These are enumerated as follows: for every nonempty row subset
$R$, form its common-neighbor set $C$, retain it when $C\ne\varnothing$
and $R$ is exactly the common-neighbor set of $C$. This is precisely the
maximality condition. Delete labels absent from the union of its triangle
supports and deduplicate by the remaining cyclic word and the two vertex
subsets. This loses no maximizing weight or response.

The exhaustive scan gives 79494 orders with winning pairs, 121 distinct
nonempty pair graphs, and 190 distinct charts. Fifty charts use one closure
vertex in each factor. Seventy use two q-vertices and one p-vertex; seventy
use one q-vertex and two p-vertices. No other chart occurs. Thus every
nonconstant weight family is an interval $\beta(z)=b+zd$, $0\leq z\leq1$.
The catalogue argument needs only the convex hull of the two vertices,
not an additional assertion about adjacency in the closure pentagon.

On each chart, solve the rank-five KKT equations exactly. The multiplier
and orbit vertices are affine in $z$, as are the radial coordinates
$\ell_i(z)$. Each radial coordinate is constant along its labeled segment.
The response therefore has the form
\[
 R(z)=\sum_i(b_i+zd_i)\max(0,h_i+zk_i).
\]
Partition $[0,1]$ at every exact hinge zero. On each resulting interval
$R$ is quadratic. Its maximum is attained at an endpoint or at the interior
stationary point of a strictly concave quadratic. The checker evaluates all
these candidates with exact field arithmetic and isolated signs for the
root $0<\tau<1$ of $\tau^4-10\tau^2+5=0$.
Of the 190 chart maxima, 110 equal zero and 80 equal $1/125$.
For example, the chart with word $(0,6,5,2,1,8,4)$ attains $1/125$ at
the parameter $z=3/10-\tau^2/50$ in the certificate's endpoint ordering.
This proves~\eqref{eq:ha-response}.
\end{proof}

The portable replay contains the literal normal construction, all positive
triangle weights, the complete cyclic-order enumeration, and the exact
piecewise-quadratic calculation. It requires Python's standard library and
a C++17 compiler. Its checker recomputes all outputs in a fresh output
directory and compares them with the retained certificate. A separate
Sage-field and Python enumeration reproduced the complete 190-chart set;
an independent replay also reproduced all exact response maxima. These
checks support the stated finite predicates; the coverage implication is
the argument above.

\subsection{Volume and the strict derivative}

\begin{proof}[Proof of Theorem~\ref{thm:ha-ascent}]
Each labeled facet has normalized cone-volume mass $1/10$. Its profile is
independent of the edge coordinate and cuts a cap from the other pentagon.
The cap depth is $\delta=2/5<1-\cos(2\pi/5)$, so the cap is a triangle
with area $\delta^2\cot(\pi/5)$. The hinge's mean on this triangle is
$\delta/3$, and the pentagon area is $A=5\tau/(1+\tau^2)$. All ten facet
averages agree. Consequently
\begin{equation}\label{eq:ha-volume}
 \int f\,d\nu=\frac{\delta^3}{3\tau A}
             =\frac8{625}-\frac{8\tau^2}{9375}.
\end{equation}
Lemma~\ref{lem:ha-repeat} and Proposition~\ref{prop:ha-certificate} identify
the supremum over all finite maximizing words, and hence their weak closure,
with $1/125$. The full-rank support hypothesis of
Theorem~\ref{thm:mc-derivative} holds because both planar normal sets have no
antipodal pairs. That theorem applies directly to the fixed twenty-row bank:
\[
 D^+\log\operatorname{sys}(H_t)
 =4\left(\frac8{625}-\frac{8\tau^2}{9375}-\frac1{125}\right)
 =\frac{20+64\sqrt5}{9375}>0.
\]
Since $H_t\to H$, the definition of a positive right derivative gives the
claimed strict improvements arbitrarily close to $H$.
\end{proof}
